\section{Nghiên cứu liên quan}
\label{sec:related_work}

Các nghiên cứu phát hiện vết dầu loang từ ảnh viễn thám SAR trong hơn ba thập kỷ qua tập trung vào ba chủ đề kỹ thuật cốt lõi: khử nhiễu đốm SAR, phân đoạn tách vùng tối ứng viên, và trích xuất đặc trưng kết hợp phân loại nhằm phân biệt vệt dầu với các vùng tối nền biển~\cite{fingas2014review,brekke2005oil}.

\subsection{Khử nhiễu đốm trên ảnh SAR}
Nhiễu đốm xuất hiện do hiện tượng giao thoa sóng tán xạ kết hợp từ vô số mục tiêu tán xạ thứ cấp phân bố ngẫu nhiên trong mỗi ô phân giải (resolution cell). Do đó, nhiễu đốm là nhiễu nhân chứ không phải nhiễu cộng: $I(x, y) = R(x, y) \cdot u(x, y)$, trong đó $R$ là độ phản xạ thực và $u$ là thành phần nhiễu ngẫu nhiên tuân theo phân phối Gamma~\cite{lee1980digital}. Các bộ lọc không gian tuyến tính kinh điển như lọc trung bình (Mean filter) hay lọc trung vị (Median filter) thường làm mờ đường biên của vệt dầu và xóa nhòa các dải hẹp~\cite{fingas2014review}.

Để khắc phục nhược điểm này, các bộ lọc thích nghi dựa trên thống kê cục bộ được phát triển, tiêu biểu là bộ lọc sai số toàn phương trung bình tối thiểu (MMSE) của Lee~\cite{lee1980digital}, bộ lọc Frost~\cite{frost1982model}, bộ lọc Kuan~\cite{kuan1985adaptive} và bộ lọc Gamma-MAP~\cite{lopes1990structure}. Bộ lọc Lee tính toán trọng số làm mịn thích nghi dựa trên tỷ số giữa phương sai cục bộ và mức nhiễu ước lượng: tại các vùng biển đồng nhất (phương sai thấp), bộ lọc tăng cường làm mịn để triệt tiêu đốm; ngược lại, tại các vùng biên ranh giới dầu--biển (phương sai cao), trọng số tiến dần về 1 để giữ lại giá trị ảnh gốc. Tuy nhiên, trên các sản phẩm SAR đã qua tiền xử lý đa góc nhìn (multi-look) hoặc chuẩn hóa mức xám về định dạng 8-bit, thống kê nhiễu nhân lý thuyết có thể bị thay đổi, khiến việc áp dụng cố định tham số số looks $L$ cần được kiểm chứng thực nghiệm cẩn trọng~\cite{brekke2005oil}.

\subsection{Phân ngưỡng và tách vùng tối ứng viên}
Sau khi làm mượt nhiễu đốm, bước tiếp theo là phân đoạn các vùng tán xạ thấp. Phương pháp phân ngưỡng toàn cục Otsu~\cite{otsu1979threshold} là kỹ thuật kinh điển nhờ tốc độ tính toán nhanh trên biểu đồ mức xám (histogram) và khả năng tối ưu hóa phương sai giữa hai lớp. Tuy nhiên, Otsu ngầm định histogram có dạng hai đỉnh (bimodal), dẫn đến nguy cơ chia đôi ảnh sai lệch khi tỷ lệ vết dầu chiếm phần quá nhỏ ($< 5\%$, chế độ $R_1$) hoặc khi khung ảnh ngập trọn vệt dầu lớn ($> 95\%$, chế độ $R_4$)~\cite{solberg2007oil}.

Để thích ứng với sự thay đổi cường độ nền biển trên diện rộng, các thuật toán phân ngưỡng cục bộ thích nghi như Niblack~\cite{niblack1986introduction} và Sauvola~\cite{sauvola2000adaptive} tính toán ngưỡng động cho từng điểm ảnh dựa trên trung bình và độ lệch chuẩn cục bộ trong cửa sổ trượt lân cận. Dù vậy, tại các vùng biển đồng nhất, ngưỡng cục bộ dễ bị hạ thấp và kích hoạt các điểm ảnh ở đuôi phân phối nhiễu thành các đốm tối giả mạo. Do đó, việc kết hợp hậu xử lý hình thái học (như lọc diện tích thành phần liên thông $A_{\min}$) và gán nhãn liên thông (Connected Component Labeling) là bắt buộc để hình thành các đối tượng vùng ứng viên có ý nghĩa vật lý trước khi phân loại~\cite{topouzelis2008oil}.

\subsection{Trích xuất đặc trưng và Phân loại vùng}
Để phân biệt giữa vệt dầu thật và các vùng tối giả mạo trên biển sinh ra từ hiện tượng biển lặng gió hay màng hữu cơ tự nhiên, các nghiên cứu nền tảng của Fiscella và cộng sự~\cite{fiscella2000oil}, Brekke \& Solberg~\cite{brekke2005oil}, Solberg và cộng sự~\cite{solberg2007oil}, Topouzelis~\cite{topouzelis2008oil} và Singha và cộng sự~\cite{singha2013oil} đã đề xuất không gian đặc trưng đa chiều trích xuất trên từng vùng ứng viên. Các đặc trưng này kết hợp cả tính chất hình học (diện tích, độ thuôn dài, độ đặc, độ phức tạp đường biên) và tính chất bức xạ/biên (độ tương phản giữa vùng tối và vành đai nền biển xung quanh, gradient biên).

Về giải thuật phân loại, các mô hình học máy truyền thống như Cây quyết định CART~\cite{breiman1984classification} và Máy học vectơ hỗ trợ (SVM)~\cite{cortes1995support} được sử dụng rộng rãi nhờ tính minh bạch, tốc độ suy luận nhanh và khả năng kiểm soát quá khớp tốt trên không gian đặc trưng kích thước nhỏ. Gần đây, các kiến trúc học sâu như mạng giám sát ngữ cảnh và biên CBD-Net của Zhu và cộng sự~\cite{zhu2022oil} (công trình giới thiệu bộ dữ liệu SOS) hay mạng phân đoạn của Krestenitis và cộng sự~\cite{krestenitis2019oil} đã đạt được độ chính xác phân đoạn điểm ảnh ấn tượng. Tuy nhiên, các mạng nơ-ron sâu đòi hỏi tài nguyên tính toán lớn (GPU chuyên dụng) và khó diễn giải tường minh quy tắc quyết định. Một quy trình tuần tự kết hợp xử lý ảnh thống kê và học máy phân loại theo vùng, được cài đặt và kiểm chứng số học chặt chẽ bằng thư viện mảng NumPy, cung cấp một đường cơ sở (baseline) nhẹ, rõ ràng và có thể kiểm chứng độc lập cho bài toán giám sát dầu tràn viễn thám.
